\section{Research question and precise scope}
\label{sec:scope}

The maintained project already has substantial exact infrastructure:
normal-coordinate validation, compressed component and disc observers,
independently checked diagram exteriors, coherent cocycle searches,
complete enumeration inside supplied sectors of small dimension,
bounded Pachner searches, and compressed cut geometry. Its most recent
source also contains an all-dimensional source-parallelity construction
with a linear bound on the inherited complementary core
\cite{repo}. The question addressed here is how to discover useful normal
disc vectors without guessing all quadrilateral types or expanding their
possibly enormous integer coordinates.

The 2021 lecture announcement is the motivation for the quasi-polynomial
target \cite{slides}. The July 2026 paper gives a new hierarchy algorithm
and polynomially checkable hierarchy certificates. Its Section 9 explicitly
leaves some step costs unspecified; Proposition 9.1 bounds iterations by
$L(g+1)^L$, and the refined construction controls hierarchy length linearly
\cite{lackenby2026}. These facts do not supply a quasi-polynomial bound for
the maintained implementation. Our results are a separate, parameterized
route to discovery.

\subsection{What is established}

The main results are the following.
\begin{enumerate}
\item A complete exact search for standard vertex meridian discs with
at most $k$ cancellation edges, using spaces of dimension at most $k+1$
on one-vertex, rank-one sources (Theorem~\ref{thm:discovery}).
\item A sparse-incidence enumeration of the determining equations with
an explicit $t^{k+O(1)}$ system-count bound
(Theorem~\ref{thm:incidence}). This is a stronger quantitative statement
than multiplying the number of edge subsets by the number of equation
subsets.
\item An extension with parameter $p+k$, where
$p=\dim H^1(M;\R)+v-1$ is the dimension of the cocycle space
(Theorem~\ref{thm:multivertex}).
\item A new explicit disc and an infinite layered family with cancellation
parameter exactly one; and, in a different class allowing multiple
vertices, a linear lower bound for every meridian.
\item Recovery of the cyclic boundary attachments of compressed
parallelity bases, together with the appropriate relative handle
descriptions for disc bases.
\end{enumerate}

None of the first three statements assumes that a useful quadrilateral
sector has already been supplied. None proves a logarithmic bound on the
new parameter for a universally available source family. The last
distinction is the remaining geometric issue, not a hidden linear-algebra
cost.

\subsection{Meaning of novelty and validation}

Normal edge coordinates, cocycle-dual surfaces, vertex normal surfaces,
and hyperplane-based ray enumeration are established tools. Edge-weight
coordinates are explicitly exposed by Regina \cite{regina}. The local
four-point calculation below is included with proof, rather than claimed
as a new fact about normal surfaces. The proposed contribution is their
combination into a complete cancellation-parameter search, its sparse
enumeration bound, its explicit source families, and the added compressed
interface operator. A literature check and comparison with the pinned
repository found no earlier statement of this particular package of
results. This is not a claim of exhaustive priority verification.

Computations are used to discover examples and test implementations.
Finite tests do not establish asymptotic complexity or replace the proofs.
Likewise, a positive certificate for a supplied triangulation does not
certify that it is the exterior of a particular input diagram. That source
correspondence remains an explicit input to any recognition-level use.

\section{Normal surfaces as compatible four-point metrics}
\label{sec:weights}

Let $\cT$ be a finite triangulation of a compact orientable 3-manifold.
Self-identifications between distinct faces are permitted, but local edge
occurrences and their orientations are retained. Write $t$ for the number
of tetrahedra and $m$ for the number of global edges; always $m\le6t$.
For a normal surface $S$, let
\[
               w_e(S)=|S\cap e|\quad(e\in E(\cT))
\]
be its geometric edge weights. Coordinates count intersections, not signed
intersection numbers and not the sum over repeated local occurrences of
a global edge.

In one tetrahedron label the vertices $0,1,2,3$. Write $t_i$ for the four
triangle coordinates and $q_{ij\mid kl}$ for the quadrilateral separating
$\{i,j\}$ from $\{k,l\}$. This quadrilateral avoids the opposite edges
$ij,kl$. At most one of the three quadrilateral coordinates is positive.
Define the three opposite-edge sums
\[
 s_0=w_{01}+w_{23},\qquad
 s_1=w_{02}+w_{13},\qquad
 s_2=w_{03}+w_{12}.
\]
For every face and every one of its corners put
\[
                 a_i^{jk}(w)=\frac{w_{ij}+w_{ik}-w_{jk}}2.
                 \tag{1}\label{eq:arcs}
\]
These are the numbers of normal arcs cutting off the corresponding corner.

\begin{lemma}[Local reconstruction]
\label{lem:local}
Nonnegative real edge weights in a tetrahedron come from a unique
admissible real standard normal vector if and only if all twelve corner
arc values in \eqref{eq:arcs} are nonnegative and the maximum of
$s_0,s_1,s_2$ is attained at least twice. If all three sums agree, all
quadrilateral coordinates are zero. Otherwise the strictly smaller sum
identifies the avoided pair of edges, and its quadrilateral coordinate is
half the difference between the larger and smaller sums. In either case
\[
                    t_i=\min_{\{j,k\}\subset\{0,1,2,3\}\setminus\{i\}}
                              a_i^{jk}(w).                 \tag{2}\label{eq:triangles}
\]
For integral edge weights, the reconstructed coordinates are integral
exactly when every face has even total edge weight.
\end{lemma}
\begin{proof}
If $q=q_{01\mid23}$ is the occupied type and $T=t_0+t_1+t_2+t_3$,
the edge equations give
\[
                    s_0=T,\qquad s_1=s_2=T+2q.
\]
The other types give the other two permutations of this relation. At a
fixed vertex, two of the corner arc values equal $t_i$, and the third
equals $t_i+q$. This proves necessity and the reconstruction formula.

Conversely suppose, for example, $s_0\le s_1=s_2$. Set
$q=(s_1-s_0)/2$ and use (2) for the triangles. Direct substitution into
the six edge equations gives
\[
 w_{01}=t_0+t_1,\quad w_{23}=t_2+t_3,\quad
 w_{ij}=t_i+t_j+q\quad(ij\notin\{01,23\}).
\]
The assumed corner inequalities give $t_i\ge0$, and the sum inequality
gives $q\ge0$. The formulas determine every coordinate, proving uniqueness.
The all-equal case is the same calculation with $q=0$.

For integer weights, the parity condition makes every number in (1)
an integer. It also makes the differences of opposite sums even: for
example $s_1-s_0$ is even by subtracting the parity equations on the faces
$012$ and $123$. Thus all reconstructed coordinates are integral.
The converse follows because each normal arc contributes two endpoints
to a face.
\end{proof}

The calculation can be pictured as the distance matrix of a tree with
four labelled leaves. The pendant lengths are $t_i$, and the single
internal length is $q$. This picture includes zero lengths and coincident
leaf locations; strict positivity is not required.

\begin{proposition}[Global edge-coordinate reconstruction]
\label{prop:global}
An assignment of nonnegative real weights to the global edges of $\cT$
is an admissible real normal vector exactly when the local conditions of
Lemma~\ref{lem:local} hold in every tetrahedron. Local reconstruction then
gives its unique standard coordinates. If the weights are rational,
clearing the denominators of those coordinates gives a genuine normal
surface.
\end{proposition}
\begin{proof}
On a paired face the three global edge weights are the same on both sides,
with their local labels permuted. Formula (1) therefore gives the same
three normal arc counts. These are precisely the normal matching
equations. Conversely a normal surface restricts to each tetrahedron as
in the lemma. Unique reconstruction on every tetrahedron gives global
uniqueness. Multiplying rational coordinates by a common denominator
preserves all matching, nonnegativity and quadrilateral constraints.
\end{proof}

\subsection{One common set of hyperplanes}

Fix a compatible choice $Q$ of at most one allowed quadrilateral type in
each tetrahedron. The image $\cC_Q$ of its standard normal cone in edge
coordinates is a pointed rational polyhedral cone. The edge map is a
linear isomorphism on this cone and its linear span: express any vector
in the span as a difference of cone vectors and apply local uniqueness.
Its defining
equalities and inequalities have their boundary hyperplanes among:
\begin{enumerate}
\item the $m$ equations $w_e=0$;
\item the three equations $s_i-s_j=0$ per tetrahedron;
\item the twelve equations $2a_i^{jk}(w)=0$ per tetrahedron.
\end{enumerate}
Denote this list by $\cH$, and write
\[
                     H=m+15t\le21t.                       \tag{3}\label{eq:H}
\]
Repeated or identically zero equations may be removed. Each row has
support on at most four global edges, including after identifications;
its coefficients have uniformly bounded absolute value. To describe a
particular $\cC_Q$, one imposes the appropriate equality between the two
larger sums, the inequality for the third sum, and the corner inequalities.
This description explains why searching the common hyperplane list does
not require prior knowledge of $Q$.

\section{Signed intersections and cancellation support}
\label{sec:cancellation}

Give every global edge an orientation. For a cooriented properly embedded
normal surface $S$, let $c_e(S)$ be the algebraic intersection number with
that edge. A face contains oriented normal arcs with cancelling endpoint
signs, so its oriented edge sum is zero. Thus $c(S)$ is an integral
1-cocycle. If an edge has $p_e$ positive and $n_e$ negative intersections,
then
\[
       w_e=p_e+n_e=|p_e-n_e|+2\min(p_e,n_e)=|c_e|+2z_e,
       \qquad z_e\in\Z_{\ge0}.                             \tag{4}\label{eq:defect}
\]

\begin{definition}
The cancellation support of the cooriented surface is
$F(S)=\{e:z_e>0\}$. Its size counts edges with cancellation, not the
number of cancelling pairs. Cancellation multiplicities may be very large.
For a triangulated solid torus, $\kappaV(\cT)$ is the minimum of
$|F(D)|$ over standard vertex normal meridian discs $D$.
\end{definition}

By a standard vertex surface we mean that its coordinate ray is an
extreme ray of the full nonnegative standard matching cone and its
coordinates are primitive integers. For an admissible vector this is
equivalent to extremality in any containing admissible coordinate sector:
such a sector is a coordinate face of the matching cone. The usual
vertex-disc existence theorem makes the set in the definition nonempty
for a triangulated solid torus \cite{hlp}. We keep the word \emph{vertex}
in the parameter because the discovery proof below uses extremality.
It does not show that an arbitrary small-cancellation disc can be replaced
by a vertex disc with equally small cancellation.

\begin{lemma}[One vertex fixes the cocycle]
\label{lem:cocycle}
Suppose $M$ is connected, orientable, and has torus boundary, $\cT$ has
one global vertex, and $b_1(M)=1$. Let $c$ be a primitive
integer generator of its cocycle lattice. An oriented essential disc with
boundary on the torus boundary has intersection cocycle $c$ or $-c$.
Consequently its weight vector satisfies
\[
                w(D)=b+2z,\qquad b=(|c_e|)_{e\in E}.       \tag{5}\label{eq:base}
\]
\end{lemma}
\begin{proof}
All 0-coboundaries vanish when there is one global vertex. The cocycle
lattice is therefore $H^1(M;\Z)$, a free abelian group of rank one.
An essential simple closed curve on a torus is nonseparating. A small
boundary collar therefore supplies a path from one side of the disc to
the other without recrossing it; equivalently the properly embedded disc
is nonseparating. A closed path crossing it once shows that its dual
cohomology class is primitive. Its cocycle is consequently one of the two
generators. Equation (5) now follows from (4).
\end{proof}

The absolute vector $b$ is itself an admissible normal vector in edge
coordinates. On a tetrahedron a cocycle is the difference of four local
heights, so $b_{ij}=|h_i-h_j|$. These are collinear four-point distances
and satisfy Lemma~\ref{lem:local}. The associated surface is the familiar
coherent cocycle-dual surface. It need not be a disc.

\begin{lemma}[The coherent positive support has small rank]
\label{lem:coherent-rank}
Let $C$ be the nonzero coherent surface of an integral cocycle $c$.
The standard matching kernel restricted to its positive support has
dimension at most $p=\dim Z^1(\cT;\R)$. In particular, when $p=1$ and
$c$ is primitive, $C$ is a primitive standard vertex surface and is
connected, although its genus need not be zero.
\end{lemma}
\begin{proof}
Fix an ordering of the original cocycle heights in each tetrahedron.
The supported coordinates are its nonzero successive height gaps: a
triangle at the minimum, the middle quadrilateral, and a triangle at
the maximum. A tied height constrains its corresponding gap to zero.
For any vector $y$ in the supported matching kernel, regard these
coordinates as formal height gaps. On each local oriented edge the
formal difference equals $\operatorname{sign}(c_e)w_e(y)$, with the
local-to-global edge orientation included. This identity is linear and
does not require $y$ to be nonnegative. If $c_e=0$, all supported normal
types avoid that edge, so $w_e(y)=0$.

Matching makes these edge weights agree at all occurrences of a global
edge. The vector $a_e=\operatorname{sign}(c_e)w_e(y)$ is consequently
a global cocycle. The map $y\mapsto a$ is injective: zero $a$ implies
zero weights, and the edge map is injective on the fixed quadrilateral
support. The dimension bound follows. If $p=1$, the nonzero vector $C$
forces dimension one. A common coordinate divisor of $C$ divides every
weight $|c_e|$, so primitive $c$ makes $C$ primitive. The primitive-ray
component argument of Proposition~\ref{prop:vertex-gate} proves
connectedness.
\end{proof}

\subsection{Computing the primitive generator}

Build the integral face matrix $A$, using one signed equation per global
triangular face and summing repeated edge occurrences. Rank-one nullspace
computation and division by the coordinate gcd produce $c$. The one-vertex
and rank-one hypotheses are checked, not inferred from a filename. Each
row of $A$ has sum of absolute coefficients at most three. Cofactors of
an independent $(m-1)$-row submatrix therefore give a generator with
\[
                    \max_e|c_e|\le3^{m-1}.                \tag{6}\label{eq:cbound}
\]
Hadamard's inequality proves this bound, and primitive division only
decreases it. Exact computation has polynomial bit complexity. In
particular, exponential geometric weights do not entail exponential
input length for this stage.

\section{Complete discovery from low-dimensional weight spaces}
\label{sec:flat-search}

For an edge subset $F$ define
\[
       W_F=\spanR\bigl(\{b\}\cup\{\mathbf e_f:f\in F\}\bigr),
       \qquad d_F=\dim W_F\le |F|+1.                     \tag{7}\label{eq:WF}
\]
Here $\mathbf e_f$ is the coordinate vector of edge $f$. We deliberately
search the whole space $W_F$. Requiring a nonnegative cancellation vector
would make a smaller search, but is unnecessary for soundness: every
returned candidate is reconstructed and topologically checked. In
particular, membership in $W_F$ alone is not an orientation certificate.

\begin{lemma}[Determining an extreme ray by active equations]
\label{lem:active}
Let $W\subset\R^m$ be a $d$-dimensional linear space and let $w\ne0$
generate an extreme ray of a sector cone $\cC_Q$. If $w\in W$, then some
$d-1$ equations from $\cH$, restricted to $W$, are independent and have
common kernel $\R w$ in $W$.
\end{lemma}
\begin{proof}
The ray stays extreme in $\cC_Q\cap W$: a nonproportional nonnegative
decomposition there would also be a decomposition in $\cC_Q$.
Consider all defining equalities and all inequalities active at $w$.
If their common kernel in $W$ had dimension greater than one, choose a
direction $u$ not parallel to $w$ in that kernel. For sufficiently small
$\varepsilon>0$, both $w+\varepsilon u$ and $w-\varepsilon u$ satisfy
every inactive strict inequality as well as the active equations. They
give a nonproportional decomposition of $2w$, a contradiction. Thus the
active kernel is precisely the line $\R w$. Select an independent
$(d-1)$-row subset. Every selected hyperplane belongs to $\cH$ by (3).
The argument includes sectors and intersections that are not full
dimensional.
\end{proof}

\begin{theorem}[Complete cancellation-supported vertex-disc discovery]
\label{thm:discovery}
For a supplied finite triangulation of a compact connected orientable
3-manifold with exactly one torus boundary, one global vertex and $b_1=1$,
there is a
deterministic exact algorithm that finds an essential disc whenever a
standard vertex essential disc has at most $k$ cancellation edges.
The algorithm makes at most
\[
       \sum_{\substack{F\subseteq E\\|F|\le\min(k,m)}}\binom H{d_F-1}
       \ \le\ \sum_{j=0}^{\min(k,m)}\binom mj\binom Hj     \tag{8}\label{eq:basic-count}
\]
small homogeneous linear-system attempts.
Each candidate has polynomially bounded encoded length and is checked
in polynomial bit time. Its overall bit cost is $t^{O(k+1)}$.
If the input is only known to satisfy the source hypotheses, exhaustion
of a smaller $k$ is inconclusive for unknot recognition.
\end{theorem}
\begin{proof}
For every $F$ with $|F|\le k$, compute a basis of $W_F$. Enumerate all
$(d_F-1)$-subsets of $\cH$ and retain those of restricted rank $d_F-1$.
Find the one-dimensional rational kernel, test its two directions for
nonnegativity, and apply Lemma~\ref{lem:local} on every tetrahedron.
Clear denominators in the reconstructed standard coordinates, divide by
their gcd, and apply the polynomial vertex-disc gate of
Proposition~\ref{prop:vertex-gate}. The zero-dimensional case contributes
no ray; a
one-dimensional space needs the empty equation subset.

If $D$ is a promised vertex disc, take its cancellation set $F$. Equations
(5) and (7) put its edge ray in $W_F$. The edge isomorphism preserves its
extremality. Lemma~\ref{lem:active} ensures that one enumerated system
recovers exactly this ray. Primitive normalization then recovers $D$.
The tests cannot manufacture a positive result: they verify genuine
normal coordinates and essential disc topology in the supplied manifold.

The bound counts all systems before pruning or duplicate removal. Exact
linear algebra, reconstruction and the vertex-disc gate are polynomial
in the source and coordinate bit lengths. The bit estimate below makes
these lengths polynomial in $t,k$ and the size of $c$. This proves the
stated parameterized bound without an assumption about the run time of
the repository's present interval-orbit engine.
\end{proof}

\begin{remark}[The displayed count]
The second inequality in (8) is used for $j\le H/2$, which includes the
relevant range $j\le m\le6t$ since $H\ge m+15t$. Then
$d_F-1\le j$ and the binomial coefficients are increasing on that range.
Zero hyperplanes can be left on the list for this upper bound.
\end{remark}

\subsection{Encoded size and verification cost}

Let $B$ bound the bit length of $b_e$ and of the source indices. A basis of
$W_F$ can be selected from $b$ and the coordinate vectors. Restricted
hyperplane coefficients have $O(B+1)$ bits. For a $d$-column system,
cofactor bounds give a null vector whose coordinates have
\[
                         O(d(B+\log(d+1)))                 \tag{9}\label{eq:bits}
\]
bits. Multiplication by the weight-space basis and the half-integer local
reconstruction preserve a polynomial bound. The produced standard
coordinates may describe exponentially many discs, but their binary
encoding and all verification inputs remain polynomial in $t,k,B$.
No estimate is measured in the expanded piece count.

\subsection{A polynomial gate for the promised vertex discs}

There is a useful way to avoid a stronger-than-needed topology subroutine.
Let $M_{\!N}$ be the standard normal matching matrix, and for a nonzero
integral admissible vector $x$ let $J=\supp x$. Its submatrix
$M_{\!N,J}$ retains exactly the positive coordinate columns.

\begin{proposition}[Rank and boundary parity certify a vertex disc]
\label{prop:vertex-gate}
Suppose $\partial M$ is a torus. A primitive nonzero integral admissible
normal vector $x$ is an essential disc if all three tests hold:
\[
 \dim\ker M_{\!N,J}=1,\qquad \chi(x)=1,\qquad
 [\beta(x)]\ne0\ \text{in }H^1(\partial M;\F),
\]
where $\beta(x)_e=w_e(x)\bmod2$ on boundary edges. Every primitive
standard vertex meridian passes these tests. All tests, including
admissibility, primitivity and the boundary coboundary quotient, have
polynomial bit complexity in the encoded source and $x$.
\end{proposition}
\begin{proof}
The nullity test says exactly that the positive support of $x$ spans a
ray of the standard matching cone. If the surface were disconnected,
each component would have an integral normal vector $y_i=\lambda_i x$
with $\lambda_i>0$ and $\sum_i\lambda_i=1$. Since $\gcd(x)=1$, there
are integers $a_j$ with $\sum_j a_jx_j=1$. Thus
$\lambda_i=\sum_j a_j(y_i)_j$ is an integer. There can be only one
component.

On every boundary face, the sum of the three intersection parities is
zero: each normal arc has two endpoints. Hence $\beta$ is a cocycle.
It is the mod-two Poincar\'e dual of the boundary curve system, because
its evaluation on a transverse edge counts precisely those intersections.
The nonzero class implies that the connected surface has nonempty
boundary. The classification of compact connected surfaces now makes
$\chi=1$ equivalent to being a disc. Its single boundary curve is
essential on the torus exactly when its mod-two homology class is
nonzero: an essential simple torus curve has primitive integral slope.

Conversely, a standard vertex disc has a one-dimensional supported
kernel, Euler characteristic one, and an essential simple boundary with
nonzero mod-two class. Matching ranks are exact rational linear algebra;
Euler characteristic is a linear normal-coordinate calculation; and
testing whether $\beta$ belongs to the image of the boundary vertex
coboundary map is linear algebra over $\F$. All matrices have $O(t)$
rows and columns. Determinant bounds control the rational bit lengths.
\end{proof}

The parity must be computed \emph{after} primitive standard-coordinate
normalization. A doubled disc has zero boundary parity. If the boundary
has several vertices, a nonzero raw edge-parity vector may still be a
coboundary, so quotienting by vertex potentials is necessary.

This gate is sufficient for the entire vertex-disc discovery theorem.
It may reject additional rays that are extreme only in the chosen weight
space; the completeness promise never required those rays. The main
search prototype instead reuses the maintained source-bound disc-component
observer and its independent replay, which can accept more surfaces.
That implementation has a bound in its realized interval traces; it is
not identified here with the abstract polynomial AHT theorem \cite{aht}.
A separate exact implementation of the rank gate accompanies the package.
Positive total Euler characteristic alone is never a positive certificate:
it can be contaminated by spheres or inessential components.

\section{Sparse incidence removes one combinatorial factor}
\label{sec:incidence}

The elementary search chooses $F$ and then chooses equations. Most such
pairs cannot have the required rank. Each equation in $\cH$ involves at
most four edges, so an invertible determining system must have a matching
between its rows and the chosen defect edges. This observation provides a
different enumeration order.

\begin{theorem}[Incidence enumeration]
\label{thm:incidence}
The completeness conclusion of Theorem~\ref{thm:discovery} can be achieved
with at most
\[
  1+\sum_{j=1}^{\min(k,m)}
   \left(4^j\binom Hj+m\,4^{j-1}\binom H{j-1}\right)       \tag{10}\label{eq:incidence-count}
\]
exact system attempts, up to duplicate descriptions. Preparation and
work per system are polynomial in $t,k,B$. Thus the bit bound has the
form $t^{k+O(1)}\exp(O(k\log(k+1)))\poly(B)$, with a fixed polynomial
exponent outside $k$.
The second summand handles the exceptional case in which all nonzero
coordinates of $b$ belong to the cancellation set.
\end{theorem}
\begin{proof}
First suppose $F$ has size $j$ and $\supp b\not\subset F$. The listed
basis $b,\mathbf e_f$ is independent. A meridian has weight
$b+u$, with $\supp u\subset F$, so its coefficient of $b$ is one.
Choose $j$ independent active equations as in Lemma~\ref{lem:active}.
Write their coefficient matrix in these basis coordinates as
$(a\mid R_F)$. The square matrix $R_F$ is invertible: otherwise a
nonzero vector supported on $F$ would lie in their common kernel, whereas
that kernel is the meridian line and every nonzero point on that line
has a nonzero $b$-coefficient.

A nonzero term in the determinant of $R_F$ matches each row to a different
edge in its support. Enumerate all $j$-subsets of the global hyperplane
list and all selections of one supported edge per row. There are at most
$4^j$ selections. Reject repeated edges and singular resulting matrices.
For each surviving edge set solve
\[
                   R_F u_F=-a.                            \tag{11}\label{eq:affine-solve}
\]
Reconstruct and check $w=b+u$. The active system for every promised
meridian is among these attempts. The method may also return harmless
extra surfaces; all positive outputs receive the same topological check.

Now suppose $\supp b\subset F$. The meridian weight is supported on
$F$ itself. In the $j$-dimensional coordinate space $\R^F$, its ray is
determined by $j-1$ independent active equations. Normalize by
$\sum_{f\in F}w_f=1$. The resulting $j\times j$ matrix is invertible,
since the active kernel is the line spanned by the nonnegative meridian
vector and the normalization is
nonzero on that ray. A nonzero determinant term matches the normalization
row to some one of the $m$ edges and matches every other row to one of at
most four supported edges. Enumerating the $j-1$ hyperplanes, these
choices, and the associated normalized systems therefore includes every
such ray. This gives the second summand in (10). Cases with actual
cancellation size below $k$ are covered at their own value of $j$.

All matrices have order at most $k+1$ and polynomially bounded entries.
The same cofactor and verification analysis applies. Substituting
$H\le21t$ into (10), and bounding the finite sum, gives the displayed
complexity form. The factor depending only on $k$ absorbs integer
arithmetic and elementary combinatorial bookkeeping of each attempt.
\end{proof}

Primitive normalization may rescale an affine candidate $b+u$. Therefore
the enumerated set $F$ is a search-space description, not automatically
the actual cancellation support of a returned disc. To certify that
parameter, recompute its final edge weights and compare them with $b$.
This does not affect completeness: a promised disc occurs at its correct
affine scale before primitive normalization.

\subsection{The radius-one specialization}

For $k=1$ the nonexceptional candidates are particularly simple. A row
$\ell\in\cH$ and an edge $e\in\supp\ell$ give
\[
                     u_e=-\frac{\ell(b)}{\ell_e}.          \tag{12}\label{eq:oneedge}
\]
There are at most $4H$ row-edge incidences. Equivalent rational changes
can be deduplicated before reconstruction. The exceptional sparse rays
are coordinate rays. This avoids rebuilding the complete hyperplane list
for every possible defect edge.

Equation (12) gives a linear number of small arithmetic candidates, not
an unconditional linear-time recognizer. A candidate can affect every
tetrahedron incident to a high-degree edge, and the complete vertex-disc gate
still has its own polynomial cost. A correct implementation reports both
the number of determining systems and the actual complete query time.

\section{Several vertices and the quasi-polynomial implication}
\label{sec:multi}

With several vertices, fixing a cohomology class does not fix its cocycle.
Integral vertex potentials add coboundaries. For a connected triangulation
with $v$ vertices, put
\[
                   p=\dim Z^1(\cT;\R)=v-1+b_1(M).
\]
Let $C$ be an $m\times p$ rational basis matrix for its cocycle space.
Every signed intersection cocycle is $C\alpha$. A fixed-class version
uses a cohomology generator together with a coboundary basis; for a knot
exterior these again have total dimension $v$.

\begin{theorem}[Sign-chamber extension]
\label{thm:multivertex}
For a finite triangulation of a compact connected orientable 3-manifold
with torus boundary and $p\ge1$, standard vertex essential discs with at most $k$
cancellation edges can be discovered in
\[
                      t^{O(p+k)}\poly(B)
\]
bit time from a supplied triangulation. Here $B$ bounds the encoded
rational basis and source. For a knot exterior $p=v$.
\end{theorem}
\begin{proof}
Consider the central arrangement in $\R^p$ with equations
$(C\alpha)_e=0$, omitting rows that vanish identically. Enumerate its
full-dimensional sign chambers. Their number is at most
\[
                    2\sum_{i=0}^{p-1}\binom{m-1}{i}.      \tag{13}\label{eq:chambers}
\]
The standard recurrence for the number of regions proves this: adding
one hyperplane creates at most as many regions as its induced arrangement
has in dimension $p-1$. Degeneracies only decrease the number.
An exact incremental enumeration using rational strict-feasibility tests
has $m^{O(p)}\poly(B)$ bit cost. Strict homogeneous inequalities may be replaced
by lower bounds of one when seeking a representative of a nonempty
chamber, since a representative can be positively rescaled.

Give each identically zero row the arbitrary sign $+1$.
For a chamber sign vector $\sigma$, let $D_\sigma$ be the corresponding
diagonal sign matrix. Every point in its closure satisfies
$|C\alpha|=D_\sigma C\alpha$, including coordinates that vanish.
Every cocycle lies in at least one such closure, so no separate
enumeration of lower-dimensional sign cells is necessary.

For each chamber and each $|F|\le k$, use
\[
             W_{\sigma,F}=\operatorname{im}(D_\sigma C)
                            +\spanR\{\mathbf e_f:f\in F\}.
\]
This space has dimension at most $p+k$. Equation (4) puts every promised
disc in one of these spaces. Apply Lemma~\ref{lem:active} and the
reconstruction procedure to its at most $p+k-1$ determining hyperplanes.
Multiplying the chamber count, the edge-subset count, and the
hyperplane-subset count gives $t^{O(p+k)}$. Rational determinant bounds
again control all bit lengths. This proves completeness without
assuming a bound on the nullity of the full quadrilateral sector.
\end{proof}

\begin{corollary}[Conditional recognition theorem]
\label{cor:qp}
Suppose a procedure constructs and verifies, from every $n$-crossing knot
diagram and in $n^{O(\log n)}$ bit time, a family of at most
$n^{O(\log n)}$ exterior triangulations of polynomial encoded size.
Suppose a fixed known constant $C$ guarantees that every unknot has in
one of these sources a standard vertex meridian with
$p+k\le C\log(n+2)$. Then complete execution of the budgeted searches above gives
unknot recognition in $n^{O(\log n)}$ bit time.
\end{corollary}
\begin{proof}
Put $r(n)=\lfloor C\log(n+2)\rfloor$. Discard sources with $p>r(n)$;
on each remaining source search only up to $k=r(n)-p$. At least one
promised source remains for an unknot, and it produces a verified disc
by the coverage assumption and Theorem~\ref{thm:multivertex}.
A nontrivial knot has incompressible torus
boundary, so no verified disc is returned. Completing all searches
therefore decides the problem. The size and parameter assumptions make
each search quasi-polynomial; the allowed number of searches preserves
that bound. Their source certificates and output certificates are
included in the accounting.
\end{proof}

This corollary identifies a concrete sufficient structural theorem. The
present work proves its search component, not its coverage hypothesis.
In particular, arbitrary triangulations cannot simply be substituted for
the controlled family: Section~\ref{sec:whitehead} gives an explicit
linear-cancellation obstruction when multiple vertices are permitted.
